\Needspace{18\baselineskip}
\section{Erasure commutation and traversal of actual deep specimens}
\label{sec:specimen-traversal}
This section connects the coordinate traversal results to the actual states of
recursive deep specimens. The main geometric construction commutes general
focus erasure with simultaneous expansion. Its coordinate equation then proves
that expanding an already focussed context agrees, up to a focus-preserving
isomorphism, with filling the erased context and inserting a focus afterwards.
The resulting traversal theorems start at the original specimen states.

For this representation we also prove the non-exiting conclusion of
Lemma~39(2), and synchronized search under the settled-mode hypotheses of
Lemma~39(3). The proof uses the paper's adjacency interpretation of entering
and exiting focuses and its output-closure obstruction \cite{paper}. A returned
coordinate determines actual rooted endpoints and is remote. Packaging it as a
new allowed specimen, and extending the argument from the checked search
relation to the full UAM, remain separate obligations. No arbitrary rewrite
relation is treated as search in these results.

\Needspace{14\baselineskip}
\subsection{Erasing a copied focus}
Fix a focussed context $C:I\to[]$, focus $f$, and a filling family $H$. Assume
that the focus is replaced by a literal corolla copy and that all other fillings
are globally focus-free. Let $t$ be an internal vertex with
$f.\mathsf{target}=t$. These data give an expanded focus on $C[H]$ whose target
is the host copy of $t$.

The abbreviations \texttt{erased}, \texttt{focus}, and
\texttt{erasedExpansion} denote, respectively, erasure of $f$ from $C$, the
retained focus on $C[H]$, and erasure of that retained focus using the explicit
host copy of $t$. The restricted family $\overline H$ assigns $H(e)$ to each
non-focus edge $e$ of the erased context. Erasure changes incidences, but not
the labels of these edges, so the replacement interfaces are unchanged.
\lean{Graph/ExpansionErasure}{14}{14}{Linear.ExpansionErasure.erased}
\lean{Graph/ExpansionErasure}{15}{15}{Linear.ExpansionErasure.focus}
\lean{Graph/ExpansionErasure}{16}{18}{Linear.ExpansionErasure.erasedExpansion}
\lean{Graph/ExpansionErasure}{19}{23}{Linear.ExpansionErasure.family}

\begin{definition}[Erasure--expansion carrier bijections]\label{term:erasure-expansion}
The vertex bijection from erasure after expansion to expansion after erasure
sends a retained host vertex to its retained host copy and an inserted internal
vertex to the same vertex in its non-focus owner's filling. The edge bijection
sends a surviving expanded edge to that same edge paired with its non-focus
owner.
\end{definition}
\begin{proof}[Verification of the bijections]
The copied focus has no internal vertices. Thus the only potential obstruction
to mapping an inserted internal vertex to a non-focus owner is impossible.
The inverse vertex map reinstates the original owner and the proof that the
host target was not selected. Both composites reduce to the original element;
proof fields agree by proof irrelevance.

For edges, a surviving edge cannot belong to the focus fibre: the copy has
exactly one edge, which is precisely the deleted focus. Hence its owner differs
from the focus edge. Conversely every edge of a non-focus filling survives.
The maps retain its owner and local edge, so both inverse equations are immediate.
These arguments apply to arbitrary corolla names, not just the literal corolla
constructor's chosen singleton name.
\end{proof}
\lean{Graph/ExpansionErasure}{24}{54}{Linear.ExpansionErasure.vertices}
\lean{Graph/ExpansionErasure}{55}{66}{Linear.ExpansionErasure.edges}

\Needspace{14\baselineskip}
\subsection{Commutation of incidences and the join coordinate}
For a host vertex different from $t$, retaining it during erasure commutes with
the vertex bijection. This is proved by the internal, input, and output vertex
cases: the resulting ports agree directly.
\lean{Graph/ExpansionErasure}{67}{73}{Linear.ExpansionErasure.keep_host}
\begin{lemma}[Contraction at a host vertex]
Contracting a host vertex in $C[H]$ and applying the carrier bijection gives the
host copy of its contraction in $C$.
\end{lemma}
\begin{proof}
If the original vertex is $t$, both contractions identify it with the focus
source. The retained-vertex calculation applies to that source, which differs
from $t$ because a focus is not a self-loop. Otherwise injectivity of the host
map proves that its image is not the expanded target; both contractions retain
the vertex and the same retained-vertex calculation applies.
\end{proof}
\lean{Graph/ExpansionErasure}{74}{93}{Linear.ExpansionErasure.contract_host}
The source and target port at index $i$ of an erased non-focus edge are the
contractions of its original port at index $i$. Indeed erasure maps contraction
over the ordered incidence lists, and indexed lookup commutes with this map.
\lean{Graph/ExpansionErasure}{94}{98}{Linear.ExpansionErasure.sourcePort}
\lean{Graph/ExpansionErasure}{99}{104}{Linear.ExpansionErasure.targetPort}
\begin{lemma}[Contraction at an inserted vertex]
For every vertex of a non-focus filling, contraction after its insertion,
followed by the carrier bijection, equals insertion into the erased context.
\end{lemma}
\begin{proof}
An inserted internal vertex is disjoint from the deleted host target and keeps
its original non-focus owner. A filling input is glued to an original source
port; apply contraction at host vertices and the source-port calculation.
A filling output is treated by the target-port calculation. These three cases
exhaust the filling's vertex type, including fillings with internal structure.
\end{proof}
\lean{Graph/ExpansionErasure}{105}{122}{Linear.ExpansionErasure.contract_insert}

\begin{theorem}[Erasure commutes with expansion]
There is a literal-label graph isomorphism
\[
  \bigl|C[H]\bigr|\ \cong\ |C|[\overline H],
\]
where the left erasure removes the retained copied focus and the right erasure
removes the original focus. Both use the specified internal target.
\end{theorem}
\begin{proof}
Use the carrier bijections above. Host and inserted vertex labels are retained;
the impossible focus-fibre internal case is excluded as before. Edge labels
are identical. For each source list, compose its insertion map, contraction,
and the vertex bijection. The inserted-vertex lemma identifies this composition
pointwise with insertion in the erased context. Equality of mapped lists proves
the source equation. The target equation is the same argument on the ordered
target list.
\end{proof}
\lean{Graph/ExpansionErasure}{123}{160}{Linear.ExpansionErasure.iso}
The isomorphism also maps the joined focus coordinate to the host copy of the
original joined coordinate. A join is the retained source of the deleted focus,
so this is exactly the retained-host-vertex calculation.
\lean{Graph/ExpansionErasure}{161}{169}{Linear.ExpansionErasure.position}
The canonical underlying graph chooses an internal target rather than taking
one as an argument. Its chosen target is necessarily the host copy of the
original chosen target, since both are equal to the expanded focus target.
Eliminating this equality transfers the isomorphism and its coordinate equation
to the canonical underlying graphs. The theorem returns the isomorphism and
the position equation together, ensuring that the same map establishes both.
\lean{Graph/ExpansionErasure}{170}{199}{Linear.ExpansionErasure.underlying_exists}

\Needspace{14\baselineskip}
\subsection{Recovering expansion of an already focussed context}
\begin{definition}[Focussed-context coordinate]\label{term:context-coordinate}
For a focussed context $f$, its coordinate in the erased graph consists of its
original mode and the joined term vertex. The term-typing proof comes from the
focus's typed singleton source. Unlike the earlier state-coordinate operation,
this definition does not require the context to be hole-free.
\end{definition}
\lean{Graph/ExpansionErasure}{200}{212}{Linear.Focussed.coordinate}
The abbreviation \texttt{erasedFamily} uses the canonical chosen target to
restrict the original family. The expanded erased graph is focus-free, since
it is isomorphic to the underlying graph of the retained expanded focus. If
the original expansion is hole-free, so is this expanded erased graph: erasure
preserves hole freedom and the isomorphism transports it. These proofs provide
the certificates required to realize coordinates as actual states.
\lean{Graph/ExpansionErasure}{213}{215}{Linear.ExpansionErasure.erasedFamily}
\lean{Graph/ExpansionErasure}{216}{220}{Linear.ExpansionErasure.expanded_free}
\lean{Graph/ExpansionErasure}{221}{228}{Linear.ExpansionErasure.expanded_closed}
\begin{theorem}[The focussed expansion bridge]
Fill the erased context, transport its original coordinate through that
filling, and realize the coordinate by focus insertion. The resulting state is
focus-preservingly isomorphic to expansion of the original focussed context.
\end{theorem}
\begin{proof}
Choose the erasure-commutation isomorphism together with its coordinate
equation. The extracted coordinate of the expanded state maps to exactly the
transported original coordinate; the modes agree as well. Realization respects
renaming, so their inserted states are focus-preservingly isomorphic. Compose
the inverse of this realization isomorphism with the general recovery
isomorphism from insertion after erasure to the original expanded state.
This composition identifies the distinguished edge and its mode, not merely
the two underlying graphs.
\end{proof}
\lean{Graph/ExpansionErasure}{229}{246}{Linear.ExpansionErasure.realization_iso}

\Needspace{14\baselineskip}
\subsection{Actual-state search and exact executions}
Removing the focus restricts every recursive edge certificate to the surviving
edges. A retained certificate keeps the same copy data; a filled certificate
keeps its hole label and related pair; a boxed certificate keeps its recursive
inner filling. None of these certificates depends on the host's changed source
or target lists, except through unchanged edge interfaces.
\lean{Graph/SpecimenTraversal}{9}{26}{Linear.DeepFillingEdge.erased}
For any typed coordinate $q$ of the erased context, \texttt{stateAt} realizes
its transported coordinate using the focus-freedom and hole-freedom certificates
just proved. At the original coordinate, the focussed expansion bridge applies.
\lean{Graph/SpecimenTraversal}{27}{32}{Linear.ExpansionErasure.stateAt}
\begin{theorem}[Search successors of the original expanded state]
Suppose the restricted family preserves search and has search origins. If the
original coordinate avoids hole boundaries, an actual-state search successor
is precisely a context-coordinate successor realized in the filling, followed
by a focus-preserving endpoint isomorphism.
\end{theorem}
\begin{proof}
Use source-isomorphism invariance of state search to replace the original
expanded source by its bridged realization. The coordinate characterization of
search from a realized state produces a step in the expanded erased graph.
Reflect it through the filling, using the boundary exclusion. Search preserves
term typing, so the reflected coordinate can itself be realized. The coordinate
equality identifies this realization with the supplied target up to the supplied
isomorphism. Conversely, transport the context step, realize it, and apply the
source bridge and endpoint isomorphism.
\end{proof}
\lean{Graph/SpecimenTraversal}{33}{58}{Linear.ExpansionErasure.search_iff}
An $n$-step context execution from the original coordinate gives an $n$-step
execution from the original expanded state class to \texttt{stateAt} of its
endpoint. Transport and realize the run, then replace its initial class using
the bridge. Similarly, a context run from the initial input to the original
coordinate proves rootedness of the original expanded state: the earlier
rootedness theorem applies to its realization, and the bridge equates its class
with the original state's class.
\lean{Graph/SpecimenTraversal}{59}{74}{Linear.ExpansionErasure.steps}
\lean{Graph/SpecimenTraversal}{75}{94}{Linear.ExpansionErasure.rooted}

\Needspace{14\baselineskip}
\subsection{Entering, exiting, and remote coordinates}
\begin{definition}[Concrete entering and exiting focuses]
\label{term:concrete-entering}
A focus is entering when its target is a source of an original shallow
hole-labelled edge. It is exiting when its source is a target of such an edge.
The predicates do not themselves restrict the mode, and are defined on the
original focussed context rather than on its filling.
\end{definition}
\lean{Graph/SpecimenTraversal}{95}{97}{Linear.Focussed.Entering}
\termdef{concrete-exiting}
\lean{Graph/SpecimenTraversal}{98}{99}{Linear.Focussed.Exiting}
\begin{lemma}[Adjacency after erasure]
The joined coordinate is a source of an erased non-focus edge exactly when
the original focus target is a source of that edge. It is a target of the erased
edge exactly when the original focus source is a target of that edge.
\end{lemma}
\begin{proof}
For sources, unpack membership in the contraction-mapped source list. A source
of a non-focus edge cannot equal the focus source, by source-incidence
uniqueness. The inverse expansion of contraction therefore recovers that source;
expanding the joined coordinate gives the original focus target. Conversely,
contraction maps the target to the joined coordinate.

For targets, the corresponding target of a non-focus edge avoids the deleted
focus target, by target-incidence uniqueness. Lifting its contraction recovers
it, while lifting the joined coordinate gives the focus source. Conversely the
focus source contracts to the joined coordinate. These arguments retain the
incidence direction; they do not identify entering and exiting with each other.
\end{proof}
\lean{Graph/SpecimenTraversal}{100}{121}{Linear.Focussed.erased_source_iff}
\lean{Graph/SpecimenTraversal}{122}{135}{Linear.Focussed.erased_target_iff}
Consequently the erased coordinate is at a hole boundary if and only if the
original focus is searching and entering, or backtracking and exiting. A hole
edge cannot be the focus edge, since their primitive labels differ; thus every
original hole witness supplies a surviving erased edge. The two adjacency
lemmas then prove both directions of the equivalence.
\lean{Graph/SpecimenTraversal}{136}{156}{Linear.Focussed.boundary_iff}

\begin{definition}[Remote coordinate]\label{term:remote-coordinate}
A coordinate is remote if its mode is not rewrite, or its vertex is not a source
of any hole-labelled edge. This is the coordinate counterpart of excluding only
an entering rewrite focus.
\end{definition}
\lean{Graph/SpecimenTraversal}{157}{161}{Linear.RemotePosition}
\begin{lemma}[Search has remote targets]
Every coordinate search successor is remote.
\end{lemma}
\begin{proof}
Search can produce rewrite mode only at a variable source or an operation source.
Such a source cannot also belong to a hole: source uniqueness would identify the
variable or operation edge with that hole edge, contradicting their labels.
An entry with arguments or a return with a next argument instead produces search
mode and is immediately remote. Nullary entry and final-argument return use the
operation-source argument, regardless of the operation's finish mode.
\end{proof}
\lean{Graph/SpecimenTraversal}{162}{193}{Linear.SearchStep.remote_target}

\Needspace{14\baselineskip}
\subsection{Specialization to recursive deep specimens}
Let $S$ be a recursive deep specimen. Its edge certificates restrict to the
erased context by the erasure theorem. Define $S.\mathsf{leftAt}(q)$ and
$S.\mathsf{rightAt}(q)$ by the corresponding \texttt{stateAt} constructions.
At the original coordinate these are focus-preservingly isomorphic to the
existing $S.\mathsf{leftState}$ and $S.\mathsf{rightState}$; the original
specimen representation has not been replaced by an assumed locality law.
\lean{Graph/SpecimenTraversal}{194}{199}{Linear.DeepSpecimen.erased_edges}
\lean{Graph/SpecimenTraversal}{200}{203}{Linear.DeepSpecimen.leftAt}
\lean{Graph/SpecimenTraversal}{204}{209}{Linear.DeepSpecimen.rightAt}
\lean{Graph/SpecimenTraversal}{210}{215}{Linear.DeepSpecimen.at_original}
The recursive certificates give search preservation and search origins on both
restricted families. Applying the actual-state search characterization yields
left and right successor equivalences whenever the original coordinate avoids
hole boundaries.
\lean{Graph/SpecimenTraversal}{216}{227}{Linear.DeepSpecimen.left_search_iff}
\lean{Graph/SpecimenTraversal}{228}{240}{Linear.DeepSpecimen.right_search_iff}
The equivalences also hold for recursive state search, with recursive
focus-preserving endpoint isomorphisms. In the forward direction, factor a
recursive search step into a literal search step and a recursive endpoint
renaming. Extract the context coordinate from the literal step and compose the
isomorphisms. In the reverse direction, take the canonical literal realization
as the intermediate state and apply the supplied recursive renaming.
\lean{Graph/SpecimenTraversal}{241}{255}{Linear.DeepSpecimen.left_deepSearch_iff}
\lean{Graph/SpecimenTraversal}{256}{272}{Linear.DeepSpecimen.right_deepSearch_iff}
In particular, a left literal search step determines a context coordinate,
a literal isomorphism from its left realization to the supplied target, and a
matching right search step to its right realization. The right step uses the
same context-coordinate witness and the identity endpoint isomorphism.
\lean{Graph/SpecimenTraversal}{273}{283}{Linear.DeepSpecimen.match_left_step}
Exact context runs give paired executions starting from the original specimen
states, by applying the general execution bridge separately on both sides.
A context run from the initial input to the original coordinate similarly proves
rootedness of both original states, by the general rootedness bridge.
\lean{Graph/SpecimenTraversal}{284}{295}{Linear.DeepSpecimen.context_steps}
\lean{Graph/SpecimenTraversal}{296}{306}{Linear.DeepSpecimen.rooted_of_context_run}

\Needspace{14\baselineskip}
\subsection{Output closure and settled traversal}
\begin{theorem}[Non-exiting rooted specimens: concrete Lemma 39(2)]
If the outer pre-template is output-closed and both states of $S$ are rooted,
then its original focus is not exiting.
\end{theorem}
\begin{proof}
An exiting witness is a hole edge with the focus source among its targets.
The singleton path consisting of this edge ends at the focus source and
contains a non-operation edge, because a hole is not an operation. The proved
recursive-specimen form of Lemma~26(3) says that one of the actual states is
not rooted, contradicting joint rootedness. The argument concerns the current
recursive representation; its finite tuple compatibility is now proved in
Section~\ref{sec:tuple-safety}. The separate comparison with set-based
hypernets remains a representation obligation.
\end{proof}
\lean{Graph/SpecimenTraversal}{307}{318}{Linear.DeepSpecimen.not_exiting}
If the focus is backtracking, or is searching without entering, then joint
rootedness and output closure exclude the coordinate hole boundary. Use the
boundary equivalence: its entering/search alternative contradicts the mode or
the non-entering premise; its exiting/backtrack alternative contradicts the
non-exiting theorem. The one-sided matching theorem therefore applies under
these settled hypotheses.
\lean{Graph/SpecimenTraversal}{319}{331}{Linear.DeepSpecimen.settled_away}
\lean{Graph/SpecimenTraversal}{332}{343}{Linear.DeepSpecimen.settled_match_left}
\begin{theorem}[Synchronized settled traversal]
Assume output closure, joint rootedness, and a backtracking or non-entering
search focus. If recursive search is enabled on either actual specimen side,
there is a typed successor coordinate $q$ such that both actual sides search to
their respective realizations at $q$. Both endpoints are rooted and $q$ is remote.
\end{theorem}
\begin{proof}
The settled hypotheses give boundary exclusion. Extract the context step from
whichever side is enabled, using its recursive-search characterization. Apply
both successor characterizations to that same step, choosing identity endpoint
isomorphisms. The search-target lemma proves remoteness of $q$. Concrete search
preserves rootedness, so each of the two constructed endpoints is rooted.
\end{proof}
\lean{Graph/SpecimenTraversal}{344}{368}{Linear.DeepSpecimen.settled_traversal}
This proves synchronized search and the required endpoint facts. It does not
yet return a new element of the original allowed specimen type: that requires
repackaging the coordinate, preserving unique holes, and transporting allowed
context membership under the relevant graph isomorphisms. The full-UAM
transition relation and its connection to these search steps also remain open.
These are the remaining interface obligations, rather than an assumed failure
of the graph-level traversal proof.

\Needspace{14\baselineskip}
\subsection{Traversal and rootedness regressions}
Return to the original deep specimen with its focus immediately before the
consumer and its sole hole inside a box. It has no shallow hole-labelled edge,
so neither entering nor exiting is possible. The boundary equivalence therefore
excludes the erased-coordinate hole boundary.
\lean{Graph/SpecimenTraversalExamples}{10}{14}{Linear.SpecimenTraversalExamples.no_boundaries}
\lean{Graph/SpecimenTraversalExamples}{15}{21}{Linear.SpecimenTraversalExamples.away}
The consumer is an actual operation of the erased context. Its singleton source
is the joined coordinate: in the focussed context its source is the focus target,
and erasure contracts that target to the join. Its argument and thunk lists
are empty. Define the returned coordinate by changing only the mode to
backtrack. With passive operation kinds, nullary entry supplies the corresponding
context step.
\lean{Graph/SpecimenTraversalExamples}{22}{42}{Linear.SpecimenTraversalExamples.consumer}
\lean{Graph/SpecimenTraversalExamples}{43}{45}{Linear.SpecimenTraversalExamples.returned}
\lean{Graph/SpecimenTraversalExamples}{46}{49}{Linear.SpecimenTraversalExamples.context_step}
The left actual-state search characterization realizes this as a step from the
original expanded specimen state. If any recursively renamed left successor is
supplied instead, extract its context coordinate. Coordinate determinism equates
it with the returned coordinate, giving its recursive endpoint isomorphism and
the matching right transition.
\lean{Graph/SpecimenTraversalExamples}{50}{53}{Linear.SpecimenTraversalExamples.actual_step}
\lean{Graph/SpecimenTraversalExamples}{54}{65}{Linear.SpecimenTraversalExamples.matching}
Finally no context run from the initial input can reach this specimen's original
coordinate. Such a run would make both original states rooted by the new bridge,
contradicting the earlier deep-box obstruction. Thus the new reconstruction and
traversal results preserve the previous distinction between well-formed states
and rooted states; a valid local search step does not imply rootedness.
\lean{Graph/SpecimenTraversalExamples}{66}{74}{Linear.SpecimenTraversalExamples.no_initial_context_run}
